% ---------------------------------------------------------------------------
% zero_polymer: everything done so far on the 566-residue reference.
%
% Code, safety checks, calibration, multi-site analysis, temperature sweeps,
% and the cost of a move (time and GPU memory). Companion to informedness.tex
% and gradient_method.tex, and to softmax/DEVLOG.txt (the chronological
% record, entries 2026-10-02 to 2026-10-15). Every number below is taken from
% a file committed in the repository; the file is named next to it.
% Numbers that are NOT in the repository are marked with \missing{...} and
% collected in Section 8 together with how to measure them.
%
% Build:  cd softmax/docs && pdflatex zero_polymer_recap.tex && pdflatex zero_polymer_recap.tex
% ---------------------------------------------------------------------------
\documentclass[11pt,a4paper]{article}

\usepackage[T1]{fontenc}
\usepackage[utf8]{inputenc}
\usepackage[margin=2.5cm]{geometry}
\usepackage{amsmath,amssymb}
\usepackage{booktabs}
\usepackage{array}
\usepackage{xcolor}
\usepackage[hidelinks]{hyperref}
\sloppy

\newcommand{\dt}{\Delta t}
\newcommand{\gradU}{\nabla U}
\newcommand{\Mass}{M}
\newcommand{\Tsftm}{T_{\mathrm{sftm}}}
\newcommand{\Uam}{U_{\mathrm{am}}}
\newcommand{\Hd}{H_d}
\newcommand{\file}[1]{\texttt{#1}}
% Placeholder for a number that is not recorded in the repository.
\newcommand{\missing}[1]{\fcolorbox{red}{red!8}{\textcolor{red!70!black}{\small\textbf{MISSING:} #1}}}

\title{zero\_polymer: a recap\\
       \large Code, validation, calibration, multi-site moves,\\
       temperature sweeps, and the cost of a move}
\author{\texttt{my\_extendedSampling} / \texttt{softmax}}
\date{2026--10--21}

\begin{document}
\maketitle

\begin{abstract}
\noindent
\texttt{zero\_polymer} is a 566-residue hemagglutinin HA0 precursor, about ten
times longer than protein~G (56 residues), on which the gradient-informed rate
sampler was developed. This note collects all the work done on it between
2026-10-02 and 2026-10-15. The sampler is correct at this length: every check
in the safety suite passes, including an exact Boltzmann comparison on a
400-state subproblem whose error decays as $1/n$. The energy scale is about
$19\times$ protein~G's, so the temperature and the step size had to be
re-derived (usable $T\approx 30$--$300$; $\dt$ rescaled as $\sqrt{T}$). One move
costs about $0.36$\,s on an RTX 4500 Ada (two forward+backward passes and one forward
pass through ESM3), and the cost does not depend on how many sites the move
changes, so a multi-site move is free in compute and is limited only by
acceptance. A first 1000-move temperature sweep locates the transition between
$T\approx 30$ and $300$ but is far from equilibrium. The production sweep
($3\times10^5$ moves) has not been launched: it waits on the study of how many
sites to change per move at each temperature.
\end{abstract}

\tableofcontents

% ===========================================================================
\section{Scope, sources, and what is missing}
% ===========================================================================
Sources, all in the repository:
\begin{itemize}
  \item \file{softmax/DEVLOG.txt}, entries from 2026-10-02 to 2026-10-15
        (the narrative and the reasoning);
  \item \file{softmax/code/zero\_polymer/*.txt}, output of the test suite run on
        the GPU;
  \item \file{softmax/monte\_carlo/zero\_polymer\_test\_1000\_1mut/} (the 1000-move
        sweep, logs committed) and
        \file{softmax/monte\_carlo/zero\_polymer\_test\_10\_300/} (the 1000-move
        test on the 10--300 grid, analysis output in \file{results.txt});
  \item the code itself: \file{classes/rate\_sampler.py},
        \file{monte\_carlo/main\_rate.py},
        \file{monte\_carlo/sweep\_temperature\_zero\_polymer.sh},
        \file{monte\_carlo/gpu\_scheduler.sh}, \file{code/tests/}.
\end{itemize}
All ESM3 runs were done on Marco's own GPU machine; nothing here was re-run for
this note. One table (Table~\ref{tab:hastings}) is new: it was computed for this
note from the committed sweep logs, and says so. Numbers that the repository
does not contain (the GPU model, wall-clock per move in a real sweep, GPU
utilisation, the production sweep and the pilot) are marked
\missing{like this} and listed in Section~\ref{sec:missing} with the command
that measures each one.

% ===========================================================================
\section{The system}
% ===========================================================================
\paragraph{Reference.} \texttt{ZERO\_POLYMER} in \file{code/tests/references.py}:
566 residues, only the 20 canonical amino acids. It is the HA0 precursor
\emph{with} its signal peptide, so the usual (mature) HA numbering is 16 behind
the array index; \texttt{mature\_to\_index()} converts. The same file defines
\texttt{zero\_polymer\_mature} and a cleaved form (HA1\,|\,HA2) used only for the
structure tests of Section~\ref{sec:structure}. The reference is selected with
the \texttt{REF} environment variable in the tests and with \texttt{--ref} in
\file{main\_rate.py}.

\paragraph{Energy.} Only the attention-map energy is used
($\lambda_{\mathrm{am}}=1$, $\lambda_{\mathrm{structure\_ce}}=\lambda_S=0$):
\begin{equation}
  \Uam(\sigma) = \sum_{i\le j}\bigl[\log a_{ij}(\sigma) - \log a_{ij}(\sigma_{\mathrm{ref}})\bigr]^2 ,
\end{equation}
where $a(\sigma)$ is the ESM3 attention map of sequence $\sigma$
(\file{utils/energies.py:compute\_U\_am}). $\Uam\ge 0$ and
$\Uam(\sigma_{\mathrm{ref}})=0$, so the reference is the global minimum. The sum
runs over $\sim\!1.6\times10^5$ pairs at $L=566$ against $\sim\!1.5\times10^3$ for
protein~G, which is why raw energies are never compared across the two
references.

% ===========================================================================
\section{The code}
% ===========================================================================
\subsection{The move (\file{classes/rate\_sampler.py})}
Each site carries one-hot logits $x_i\in\mathbb{R}^{K}$ ($K=25$ vocabulary
entries, of which 20 are canonical). One move of
\texttt{ExtendedProteinRateSampler.\_step}:
\begin{enumerate}
  \item Draw $n$ distinct sites uniformly (\texttt{n\_sites}, default 1). The
        choice is symmetric and cancels from the acceptance ratio.
  \item Gradient at the current sequence $A$: only the chosen sites are relaxed,
        $\mathrm{softmax}(x_i/\Tsftm)$ with $\Tsftm=0.1$; every other site stays
        exactly one-hot. One forward and one backward pass through ESM3
        (\texttt{\_grad\_pass\_sites}).
  \item One leapfrog half-step with fresh momentum $p\sim\mathcal N(0,\Mass T)$,
        \begin{equation}
          \delta x = \frac{\dt}{\Mass}\,p - \frac{\dt^2}{2\Mass}\,\gradU
          \;\sim\; \mathcal N\!\bigl(\mu,\sigma^2\bigr),\qquad
          \mu=-\frac{\dt^2}{2\Mass}\gradU,\quad \sigma=\dt\sqrt{T/\Mass}.
        \end{equation}
        The target at each site is the argmax of $\delta x$ over the 19
        competitors (canonical residues other than the current one).
  \item Forward proposal probability, per site, in closed form
        (\file{utils/proposals.py:log\_pointing\_prob}, 40 Gauss--Hermite nodes):
        \begin{equation}
          a_{i\to j}=\int\!\phi(u)\prod_{k\ne i,j}\Phi\!\bigl(\alpha_{kj}+u\bigr)\,du,
          \qquad \alpha_{kj}=\frac{\mu_j-\mu_k}{\sigma}=\frac{\dt\,(\partial_k U-\partial_j U)}{2\sqrt{\Mass T}} .
        \end{equation}
        The joint proposal is the product over sites.
  \item Substitute, take the gradient again at the candidate $B$ (second
        forward+backward), and compute the reverse probability $a_{B\to A}$.
  \item Exact energy of $B$: one forward pass on the discrete sequence, no
        gradient (\texttt{\_exact\_energy}).
  \item Accept with probability $\min(1,e^{\ell})$,
        $\ell = -\Delta U/T + \log a_{B\to A} - \log a_{A\to B}$.
\end{enumerate}
This is an exact Metropolis--Hastings kernel for $e^{-U/T}$
(informedness.tex). Cost per move: \textbf{2 forward+backward passes and 1
forward pass}, independent of $n$ (Section~\ref{sec:cost}).

Changes made during the zero\_polymer work:
\begin{itemize}
  \item \textbf{n\_sites} (2026-10-15): several sites per move, one relaxed
        gradient pass for all of them. With \texttt{n\_sites}$>1$ the
        \texttt{site}, \texttt{cur\_aa} and \texttt{proposed\_aa} columns of
        \file{data.dat} are comma-separated, and $\Delta U$ and the log
        probabilities are sums over sites. Checked by
        \file{code/tests/test\_nsite\_kernel.py} (exact detailed balance and
        stationarity on 3-letter toy systems for $n=1,2,3$, sampled targets
        against pointing probabilities, restart reproducibility). The test
        passes on the GPU machine after a float32 tolerance fix
        (commit \texttt{3a85af7}).
  \item \textbf{Canonical starts} (2026-10-15): \texttt{init\_muts} used to draw
        from all 25 vocabulary entries; a site holding X/B/U/Z/O can never be
        left (its reverse probability is 0), so 20\% of start mutations froze a
        site. \texttt{\_mutate\_canonical} now draws from the 20 canonical
        residues, and \texttt{\_setup} refuses a non-canonical reference. Runs
        with \texttt{init\_muts}$=0$ (every sweep so far) are unaffected.
  \item \textbf{Entropy column removed} (2026-10-13): it was the entropy of
        $\mathrm{softmax}(x/\Tsftm)$ over one-hot sites, i.e.\ the constant
        $L\times 0.011973=6.777$ nats. Old logs still carry it.
\end{itemize}

\subsection{Driver and sweep}
\begin{description}
  \item[\file{monte\_carlo/main\_rate.py}] reads \file{pars.txt},
        \file{settings.txt} and \file{CONSTANTS.txt}, takes \texttt{--ref}
        (previously hard-wired to protein~G), and \texttt{--extend}, which
        continues a finished run to a larger \texttt{moves} from its last
        checkpoint, appending to \file{data.dat} and recording the extension in
        \file{extensions.txt}. The extension was tested on the GPU (T=300,
        1000$\to$2000 moves) and worked; its logs are not committed.
  \item[\file{monte\_carlo/sweep\_temperature\_zero\_polymer.sh}] writes one
        \file{pars\_T<T>.txt} per temperature and launches \file{main\_rate.py}
        through the GPU scheduler. Environment overrides: \texttt{TS\_OVERRIDE},
        \texttt{MOVES}, \texttt{LOG\_STEP}, \texttt{PRINT\_STEP},
        \texttt{INIT\_MUTS}, \texttt{N\_SITES}, \texttt{DT\_AT\_T2} (63),
        \texttt{DT\_EXP} (0.5), \texttt{EXTEND}, \texttt{DRY\_RUN},
        \texttt{RUN\_DIR}, \texttt{MAX\_RUNNING} (2), \texttt{MIN\_FREE\_MB}
        (12000). The step size follows
        \begin{equation}\label{eq:dt}
          \dt(T) = 63\,\bigl(T/2\bigr)^{1/2},
        \end{equation}
        which holds $\dt/\sqrt{T}$, and therefore the informedness of the
        proposal, fixed across the sweep (Section~\ref{sec:dt}).
  \item[\file{monte\_carlo/gpu\_scheduler.sh}] gates each launch on free GPU
        memory (\texttt{nvidia-smi}) and on a process cap, and now reports each
        job's exit status (a crash used to print as ``finished'').
  \item[\file{monte\_carlo/analyze\_ensemble.py}] per-temperature analysis of
        the checkpoints: $U$ and $\Hd$ trajectories, site entropy, pairwise
        overlap $q$, and a burn-in check. It needs \texttt{--ref-seq}, since it
        defaults to protein~G.
\end{description}

\subsection{The safety-check suite (\file{code/tests/}, outputs in \file{code/zero\_polymer/})}
\begin{center}
\small
\begin{tabular}{@{}p{0.34\textwidth}p{0.58\textwidth}@{}}
\toprule
script & what it measures \\
\midrule
\file{scan\_U\_am\_landscape.py} & numerical health of $\log a$, noise floor, $U(\mathrm{ref})=0$, all single mutants from the reference \\
\file{test\_steepest\_descent.py} & does the gradient's argmax pick a good substitution (top-1, $\Delta U$ at the pick) \\
\file{compare\_grad\_methods.py} & site-only vs whole-sequence gradient \\
\file{test\_informedness\_vs\_dt.py} & $p_{\mathrm{match}}(\dt)$, the probability the noisy draw follows the gradient \\
\file{validate\_boltzmann\_toy.py} & full sampler on 2 free sites (400 states) against exact enumeration; prints a $T$ calibration table \\
\file{test\_joint\_coupling\_N.py} & coupling of $N$ simultaneous substitutions \\
\file{test\_multisite\_tradeoff.py} & cost per operation, landscape per starting point, coupling vs $N$ and drift, throughput $a_N N$ vs $a_1$ \\
\file{test\_attention\_locality.py} & how local the change in $a$ is after one substitution; gradient cost vs number of relaxed rows; backward reproducibility \\
\file{test\_nsite\_kernel.py} & exactness of the \texttt{n\_sites} kernel (toy energy, no ESM3) \\
\bottomrule
\end{tabular}
\end{center}

% ===========================================================================
\section{Validation and calibration results}
% ===========================================================================
\subsection{The energy is numerically sound}
From \file{scan\_U\_am\_results.txt}: $320\,356$ attention entries, none zero, none
non-finite, minimum $5.4\times10^{-5}$, $\log a\in[-9.82,-2.20]$. Running the
forward pass twice gives a noise floor of exactly 0: the forward is
bit-reproducible, so every $\Delta U$ is signal. $\Uam(\mathrm{ref})=0$ and no
single mutant goes below it. The control run on protein~G reproduces every
recorded protein~G number to the last decimal, so the refactored harness
(\file{references.py}) changed nothing.

\subsection{Energy scale}
Mean $\Delta U$ of a single substitution from the reference is $603.35$ on
zero\_polymer against $31.44$ on protein~G, a factor $19.2$ for a length factor
$10.1$, i.e.\ roughly $\Delta U\propto L^{1.3}$. This one number is why every
$T$ inherited from protein~G is wrong here: at the old working point $T=2$ a
typical move has $\Delta U/T\sim 10^2$ and is never accepted.

\subsection{Gradient quality}
From \file{steepest\_descent\_results.txt} and
\file{steepest\_descent\_S30\_M51.txt}:
\begin{center}
\begin{tabular}{@{}lccccc@{}}
\toprule
 & drift & top-1 & $\Delta U$ at pick & $\Delta U$ other & advantage \\
\midrule
protein G, MUTS=5     & 8.9\% & 2/10  & $+17.20$  & $+20.99$  & 3.79 \\
zero\_polymer, MUTS=5 & 0.9\% & 2/10  & $+474.38$ & $+542.46$ & 68.08 \\
zero\_polymer, MUTS=51 & 8.8\% & 10/30 & $-360.03$ & $+244.25$ & 604.28 \\
\bottomrule
\end{tabular}
\end{center}
At comparable drift the gradient is \emph{more} informative on the longer
sequence (73\% of picks downhill). The bad average at MUTS=5 came entirely from
two proline picks next to the reference ($+3872$, $+925$); in the bulk the
gradient picked proline three times, all downhill. Measuring at 0.9\% drift is a
floor measurement and misleads; the lesson is to test at a drift comparable to
the one the chain actually runs at. The site-only and whole-sequence gradients
agree (cosine $\ge0.9999$ at all 11 sites tested).

\subsection{Step size}\label{sec:dt}
The working point $p_{\mathrm{match}}\approx0.15$ sits at $\dt\approx63$ at $T=2$
on zero\_polymer, against $\dt\approx1270$ on protein~G
(\file{informedness\_vs\_dt\_results.txt}); the factor $\sim20$ matches the ratio of
gradient magnitudes. Because the proposal depends on $\dt$ only through
$\alpha\propto\dt/\sqrt{\Mass T}$, the step must be rescaled with $\sqrt{T}$ when
$T$ changes, which is Eq.~\eqref{eq:dt}: $\dt=460$ at $T=107$, $\dt=772$ at
$T=300$. These $p_{\mathrm{match}}$ values were measured at 0.9\% drift; they
have not been re-measured in the bulk.

\subsection{Temperature window (2-site subproblem)}
From \file{boltzmann\_Tcal.txt}: two sites (30, 52) free, the rest at the
reference, 400 states enumerated exactly. $U$ ranges over $[0, 6235]$, median
$740$.
\begin{center}
\begin{tabular}{@{}rrrl@{}}
\toprule
$T$ & effective states & max $p$ & verdict \\
\midrule
2.0    & 1.00   & 1.0000 & frozen (the protein~G $T$) \\
10.7   & 1.05   & 0.9743 & frozen \\
32.2   & 2.97   & 0.5475 & usable \\
107.4  & 41.66  & 0.0804 & usable \\
322.2  & 183.45 & 0.0149 & usable \\
1074   & 337.19 & 0.0050 & flat \\
\bottomrule
\end{tabular}
\end{center}
This window describes a 2-site problem next to the reference; it is a starting
bracket for the full chain, not its answer.

\subsection{Boltzmann validation}
From \file{validate\_boltzmann\_results\_dt460.txt},
\file{\dots\_dt63.txt} and \file{\dots\_dt460\_M40000.txt}, $T=107$:
\begin{center}
\begin{tabular}{@{}lccc@{}}
\toprule
 & $\dt=460$, $n=16\,000$ & $\dt=63$ (near uniform) & $\dt=460$, $n=32\,000$ \\
\midrule
acceptance          & 0.306  & 0.289  & 0.309 \\
KL(emp$\,\|\,$exact) & 0.0316 & 0.0338 & 0.0169 \\
total variation     & 0.0787 & 0.0848 & 0.0569 \\
states visited      & 224/400 & 223/400 & 252/400 \\
\bottomrule
\end{tabular}
\end{center}
Doubling $n$ divides KL by $1.87$ (expected 2 for $1/n$) and TV by $1.38$
(expected $\sqrt2=1.41$), with a stable autocorrelation time of $\sim25$ moves.
The error decays as an unbiased sampler's must, in two independent functionals.
This closes validation for the restricted problem; the full 566-site chain
cannot be checked against anything exact.

\subsection{Gradient locality and reproducibility}
From \file{attention\_locality.txt}. One substitution changes the attention map
almost everywhere: about 98\% of $|\Delta a|$ and of $\Delta\Uam$ comes from
pairs that do not involve the mutated site (per entry the change is 6--7.5$\times$
stronger near the site, but there are 282$\times$ more distant entries). At site
80 the local and distal contributions have opposite signs. Separately, the
backward pass is not bit-reproducible under bfloat16 autocast: two gradients on
the same input differ by 0.18\% and change the preferred substitution at 3/566
sites. This does not affect correctness (the kernel uses the probabilities it
actually drew), but it means a gradient diagnostic will not re-run identically.

% ===========================================================================
\section{Multiple mutations per move}
% ===========================================================================
\paragraph{Coupling} (\file{test\_joint\_coupling\_N\_results.txt}). Coupling is
$c = \Delta U_{\mathrm{joint}} - \sum\Delta U_{\mathrm{isolated}}$. Relative to the
move, coupling is 2--5$\times$ weaker on zero\_polymer than on protein~G
($|c|/|\Delta U_{\mathrm{joint}}| = 0.08$--$0.47$ against $0.43$--$0.85$) and mostly
favourable in sign. Two qualifications came later: part of this is a distance
effect, since sites 30 residues apart couple 1.4--2.6$\times$ more than random
sites $\sim190$ apart (\file{multisite\_tradeoff.txt}, section 3), and an $N$-site
move still costs about $N$ single-site moves in $\Delta U$.

\paragraph{Cost parity} is exact (Section~\ref{sec:cost}): the $N$ mutations come
for free in compute, so the design question is only whether
$a_N\,N > a_1$, where $a_N$ is the acceptance of an $N$-site move.

\paragraph{Throughput} (\file{multisite\_tradeoff.txt}, section 4). It depends on
where the chain sits. Acceptance of the median move at $T=107$:
\begin{center}
\begin{tabular}{@{}rccc@{}}
\toprule
$N$ & drift 5 & drift 51 & drift 150 \\
\midrule
1  & 0.133 & 1 & 1 \\
5  & $7.9\times10^{-8}$ & 1 & 1 \\
12 & $2.6\times10^{-15}$ & 1 & 1 \\
20 & $4.0\times10^{-22}$ & 1 & 1 \\
\bottomrule
\end{tabular}
\end{center}
Near the reference $a_N$ collapses geometrically and multi-site loses. Far from
it the median move is downhill at every $N$ and multi-site wins, but those
starting points are random sequences well above equilibrium: that is
relaxation, not sampling. (The real sampler accepts $0.31$ at $T=107$ in the
Boltzmann test, against $0.6$--$0.8$ for single-site moves from these points.)
The equilibrium question is open and is what the 10-site pilot is for:
\texttt{INIT\_MUTS=51}, $T\in\{107,300\}$, 5000 moves, \texttt{N\_SITES=10}
against \texttt{N\_SITES=1}; compare $10\,a_{10}$ with $a_1$ after move 2500.
Pilot result: \missing{pilot\_n10 / pilot\_single acceptance and $\Hd$ after
move 2500}.

% ===========================================================================
\section{Temperature sweeps}
% ===========================================================================
\subsection{\texorpdfstring{1000-move sweep, $T=0.1$ to $10^4$}{1000-move sweep, T=0.1 to 1e4}}
\file{zero\_polymer\_test\_1000\_1mut/}, single-site, start at the reference,
$\dt$ from Eq.~\eqref{eq:dt}. The first attempt lost 9 of 11 runs to CUDA
out-of-memory on move 1 (Section~\ref{sec:gpu}); the re-run (commit
\texttt{35757fa}) completed all 11. Window: moves 500--1000.
\begin{center}
\begin{tabular}{@{}rrcrrr@{}}
\toprule
$T$ & $\dt$ & accept & $\Hd$, move 500$\to$1000 & mean $\Hd$ & mean $U$ \\
\midrule
0.1    & 14  & 0.000 & 0 $\to$ 0     & 0.0   & 0 \\
1      & 45  & 0.000 & 0 $\to$ 0     & 0.0   & 0 \\
10     & 141 & 0.010 & 0 $\to$ 3     & 1.2   & 35 \\
31.6   & 250 & 0.052 & 14 $\to$ 32   & 24.0  & 859 \\
56.2   & 334 & 0.168 & 54 $\to$ 114  & 87.0  & 4023 \\
100    & 445 & 0.325 & 99 $\to$ 198  & 153.0 & 7979 \\
177    & 593 & 0.489 & 161 $\to$ 290 & 231.4 & 14361 \\
316    & 792 & 0.557 & 196 $\to$ 313 & 263.5 & 20392 \\
562    & 1056 & 0.639 & 228 $\to$ 357 & 302.5 & 29016 \\
1000   & 1409 & 0.665 & 237 $\to$ 370 & 311.9 & 33445 \\
$10^4$ & 4455 & 0.491 & 241 $\to$ 347 & 295.7 & 42648 \\
\bottomrule
\end{tabular}
\end{center}
Frozen for $T\le1$, nearly frozen at $T=10$; acceptance climbs from 5\% to 56\%
between $T\approx30$ and $300$. Nothing is equilibrated: $\Hd$ is still rising by
60--135 between moves 500 and 1000 at every $T\ge56$. This is a pipeline test,
not a measurement of the equilibrium ensemble.

\paragraph{Why acceptance drops at $T=10^4$.} The DEVLOG left this open. Splitting
$\ell$ into its energy part $-\Delta U/T$ and its proposal (Hastings) part
$h=\log a_{B\to A}-\log a_{A\to B}$ in the committed logs gives
Table~\ref{tab:hastings} (computed for this note, moves 500--1000).
\begin{table}[t]
\centering
\begin{tabular}{@{}rcccc@{}}
\toprule
$T$ & accept & fraction downhill & accept $|$ downhill & mean $h$ \\
\midrule
56.2   & 0.168 & 0.052 & 0.96 & $-0.26$ \\
100    & 0.325 & 0.138 & 0.96 & $-2.07$ \\
177    & 0.489 & 0.236 & 0.89 & $-0.08$ \\
316    & 0.557 & 0.301 & 0.85 & $-1.33$ \\
562    & 0.639 & 0.383 & 0.72 & $-0.96$ \\
1000   & 0.665 & 0.493 & 0.67 & $-3.37$ \\
$10^4$ & 0.491 & 0.615 & 0.38 & $-1.24$ \\
\bottomrule
\end{tabular}
\caption{Acceptance split by the sign of $\Delta U$, from the status rows of
\file{zero\_polymer\_test\_1000\_1mut/logs\_temperature\_zero\_polymer/T*.log}.
Computed for this note; not in the DEVLOG.}
\label{tab:hastings}
\end{table}
At high $T$ the energy term vanishes and acceptance is set by $h$ alone. A
growing share of \emph{downhill} proposals is rejected (38\% accepted at
$T=10^4$): with $\dt\propto\sqrt T$ the forward move follows the gradient
sharply, and the reverse move back to the old residue is unlikely. This
supports the DEVLOG's first suspect (the proposal at $\dt\approx4450$) but
does not prove it; a run at $T=10^4$ with a smaller $\dt$ would. It does not
matter for the working range $T\le300$.

\subsection{\texorpdfstring{1000-move test on the production grid, $T=10$--$300$}{1000-move test on the production grid, T=10-300}}
Quarter-decade grid, $\dt$ from Eq.~\eqref{eq:dt}, start at the reference.
Acceptance and $\Hd$ are from the per-move logs
(\file{logs\_temperature\_zero\_polymer/T*.log}, pushed 2026-10-22 on branch
\texttt{claude/project-thread-zyqp1u}), moves 500--1000; site entropy and $q$
are from \file{results.txt} (\file{analyze\_ensemble.py} on the 6 checkpoints
after move 500).
\begin{center}
\begin{tabular}{@{}rcrrrrc@{}}
\toprule
$T$ & accept & $\Hd$, 500$\to$1000 & $U$ at 1000 & downhill & site entropy & mean $q$ \\
\midrule
10    & 0.010 & 0 $\to$ 3     & 75     & 0.002 & 0.0040 & 0.997 \\
17.8  & 0.024 & 4 $\to$ 14    & 395    & 0.004 & 0.0120 & 0.990 \\
31.6  & 0.052 & 14 $\to$ 32   & 1226   & 0.010 & 0.0240 & 0.981 \\
56.2  & 0.168 & 54 $\to$ 114  & 5044   & 0.052 & 0.0796 & 0.936 \\
100   & 0.325 & 99 $\to$ 198  & 8970   & 0.138 & 0.1512 & 0.880 \\
177.8 & 0.489 & 161 $\to$ 291 & 16797  & 0.224 & 0.2257 & 0.821 \\
300   & 0.613 & 200 $\to$ 317 & 22059  & 0.345 & 0.2675 & 0.792 \\
\bottomrule
\end{tabular}
\end{center}
(``downhill'' is the fraction of proposals with $\Delta U<0$.) At $T=10$, 31.6,
56.2 and 100 the numbers are identical to the first sweep: same seed, same $\dt$,
same trajectory, so the runs are reproducible move for move. Same picture as
before: $U$ and $\Hd$ still rising at every $T$. The ``looks settled''
verdict that \file{analyze\_ensemble.py} prints for $T\ge56$ refers to the site
entropy over the last 500 moves and should not be read as equilibration.

\subsection{Production sweep}
Plan (DEVLOG 2026-10-14/15): the same grid, $3\times10^5$ moves per run (the
paper's minimum for protein~G, Zambon et al.\ 2024, Sects.~2.1 and 3.1),
\texttt{LOG\_STEP=PRINT\_STEP=1000}, extendable to $10^6$ with
\texttt{EXTEND=1}. Keeping the paper's $\sim5.4\times10^3$ proposals per site at
566 sites would need $3\times10^6$ moves, about 12.7 days per run, so it was
deferred. Status (2026-10-22): \textbf{not launched}. It waits on the study of
the number of sites per move (\texttt{n\_sites}) at several temperatures, which
picks the move to use before the full simulation is run (branch
\texttt{claude/project-thread-zyqp1u}).

The paper's equilibration criterion (drop the transient while the Hamming
autocorrelation to the reference exceeds 0.1, then non-overlapping block
averages) is \emph{not} implemented; \file{analyze\_ensemble.py} uses a fixed
50\% burn-in. \file{data.dat} has $\Hd$ for every move, so it can be added.

% ===========================================================================
\section{Cost of a move: time and GPU memory}\label{sec:cost}
% ===========================================================================
\subsection{Time per move}
Measured with CUDA synchronisation at $L=566$
(\file{multisite\_tradeoff.txt} section 0, \file{attention\_locality.txt}
section 2):
\begin{center}
\begin{tabular}{@{}lr@{}}
\toprule
operation & time \\
\midrule
forward pass, no gradient ($t_{\mathrm{exact}}$)   & 60.3 ms \\
forward + backward, 1 site relaxed                 & 150.4 ms \\
forward + backward, all 566 sites relaxed          & 150.0 ms \\
\quad of which backward                            & $\approx$ 89.5 ms \\
\midrule
model of one move, $2\times150.4 + 60.3$           & \textbf{361 ms} \\
model of one $N$-site move                         & 360 ms (ratio 0.998) \\
\bottomrule
\end{tabular}
\end{center}
Forward+backward time with $k$ relaxed rows, $k=1,2,5,20,100,566$:
150.8--151.1\,ms, flat (ratio $k{=}1$/$k{=}L$ = 0.9985). The reason is that
self-attention is global, so the backward has to cross all 48 layers at every
position whatever $k$ is; only the last projection, $O(LK)$, depends on $k$.

In a real run (the extension test, $T=300$, one job alone on the card,
measured 2026-10-22):
\begin{itemize}
  \item CPU time, \file{data.dat}'s \texttt{time} column
        (\texttt{time.process\_time()}) at move 1000: \textbf{0.364\,s/move};
  \item wall-clock, from the modification times of the checkpoints at move 0
        and move 1000 (364\,s apart): \textbf{0.364\,s/move}.
\end{itemize}
Both agree with each other and with the 361\,ms model, so for this code the CPU
time column is a good proxy for wall-clock (the process waits on the GPU
synchronously). (The third checkpoint, at move 2000, is 7974\,s after the
second; that gap includes the idle time between the original run and the
\texttt{--extend} launch, so it says nothing about speed.) With \textbf{two jobs} sharing the GPU (n\_sites study, $T=316$,
\texttt{n\_sites=1}, measured 2026-10-22) each job runs at \textbf{0.789\,s/move}
wall-clock (checkpoints 394\,s per 500 moves; 6308\,s for moves 0--8000; CPU
time 6311\,s). That is $2.17\times$ the single-job time, so the combined
throughput (2.53 moves/s) is about 8\% \emph{below} one job alone (2.75
moves/s): one job already saturates the card, and running two at once buys
nothing. Each process also shows 100\% CPU, so whether the limit is the GPU or
the per-move Python work on the CPU is not settled; GPU utilisation from
\texttt{nvidia-smi} would tell. The cost is the same at every $T$ (every move does the same three
passes); only acceptance changes.

Derived figures at 0.364\,s/move, one job: $\approx 9.9\times10^3$ moves per hour per job;
$3\times10^5$ moves $\approx1.3$ days per run, $10^6$ $\approx4.2$ days; the
7-point grid at two jobs at a time $\approx5$ days for $3\times10^5$, plus
$\approx12$ more to extend to $10^6$. Accepted substitutions per hour per job are
acceptance $\times\,9.9\times10^3$, e.g.\ $\approx3.2\times10^3$ at $T=100$ and
$\approx5.5\times10^3$ at $T=316$ (acceptance from the 1000-move sweep, not at
equilibrium).

For comparison, the original \file{ep\_sampler.py} runs 100 leapfrog steps (100
forward+backward passes) per Metropolis move. Its time per move on
zero\_polymer, needed for the speed-up figure, is not recorded:\\
\missing{measured time per move of \file{ep\_sampler.py} on zero\_polymer}.

\subsection{GPU memory}\label{sec:gpu}
\begin{center}
\begin{tabular}{@{}lr@{}}
\toprule
quantity & value \\
\midrule
GPU                                                      & NVIDIA RTX 4500 Ada Generation \\
card memory (total / free when idle, \texttt{nvidia-smi}) & 24\,570 / 23\,853 MiB \\
peak allocated by PyTorch, one forward+backward          & 6\,369.6 MB, flat in $k$ (+0.2 MB at $k=L$) \\
one sweep job, from \texttt{nvidia-smi} (2026-10-13)    & 4.6--8.7 GB \\
jobs run at once                                         & 2 (\texttt{MAX\_RUNNING}) \\
launch gate                                              & 12\,000 MB free (\texttt{MIN\_FREE\_MB}) \\
GPU utilisation during a sweep                           & \missing{\% busy, 1 and 2 jobs} \\
\bottomrule
\end{tabular}
\end{center}
The out-of-memory failure: the scheduler launched four jobs back to back, and
each check read the full 23\,853 MB free because the earlier jobs had not
allocated yet. Four jobs then filled the card and nine runs died on move 1. The
fix raised the gate to 12\,000 MB and capped concurrency at 2; three jobs fit at
typical use ($\sim$5 GB) but not at the 8.7 GB peak, so 2 was kept. A 10-site
move relaxes 10 rows instead of 1, which adds $\sim$0.2 MB at most by the flat-$k$
measurement, so the pilot should fit the same budget; \missing{measured peak of
an \texttt{N\_SITES=10} job}.

% ===========================================================================
\section{Structure tests on the cleaved form}\label{sec:structure}
% ===========================================================================
Not part of the sampler, recorded for completeness (DEVLOG 2026-10-08). Cleaving
at Arg329 (mature numbering) into HA1\,|\,HA2 with the chainbreak token works
through the stock ESM3 SDK. At 256 decoding steps neither the cleaved nor the
mature form folds (pTM 0.23 and 0.26, extended chain, bridge CB--CB 96 and
130\,\AA), so the chainbreak is not the cause; the likely reason is that HA is an
obligate trimer. The positive control (protein~G through the same path) has not
been run, and until it is, these numbers do not say anything about ESM3. The
sampler computes no structure, so \texttt{num\_steps} in \file{CONSTANTS.txt} is
inert for it.

% ===========================================================================
\section{Missing numbers and how to get them}\label{sec:missing}
% ===========================================================================
Run on the GPU machine; paths are relative to a sweep's \texttt{RUN\_DIR}.
\begin{enumerate}
  \item \textbf{GPU memory and utilisation during a sweep.} Start before the
        sweep, stop after, and report the peak per job and the mean
        utilisation with 1 and 2 jobs:
\begin{verbatim}
nvidia-smi --query-gpu=timestamp,utilization.gpu,memory.used \
           --format=csv -l 60 > gpu_usage.csv &
nvidia-smi --query-compute-apps=pid,used_memory \
           --format=csv -l 60 > gpu_apps.csv &
\end{verbatim}
  \item \textbf{Pilot.} Accept and mean $\Hd$ after move 2500, in each
        \file{pilot\_n10} and \file{pilot\_single} run, plus the peak memory of
        one \texttt{N\_SITES=10} job from item 2:
\begin{verbatim}
awk -F'\t' '$1 ~ /^[0-9]+$/ && $1 >= 2500 {n++; a += $14; h += $6}
            END {print "accept", a/n, "mean Hd", h/n}' data.dat
\end{verbatim}
  \item \textbf{Production sweep} ($3\times10^5$ moves, not launched yet):
        once finished, the same table as Section~6.1 over the last half of
        each run.
  \item \textbf{Baseline speed-up.} Time per move of \file{ep\_sampler.py}
        (\file{main.py}) on zero\_polymer at the same $T$, and its acceptance,
        so that accepted substitutions per hour can be compared between the two
        samplers.
\end{enumerate}

% ===========================================================================
\section{Summary}
% ===========================================================================
\begin{center}
\small
\begin{tabular}{@{}p{0.42\textwidth}p{0.5\textwidth}@{}}
\toprule
question & answer \\
\midrule
Is the sampler correct at $L=566$? & Yes, on every check; exact 2-site test passes with $1/n$ decay. \\
Energy scale vs protein G & $\Delta U$ per substitution $19\times$ larger ($\propto L^{1.3}$). \\
Usable temperature & $T\approx30$--$300$ (transition region in the 1000-move sweep). \\
Step size & $\dt=63\sqrt{T/2}$ ($p_{\mathrm{match}}\approx0.15$, measured at the floor). \\
Is the gradient useful here? & Yes in rank: 33\% top-1 and downhill on average at 8.8\% drift. \\
Time per move & $\approx0.36$ s (2 fwd+bwd + 1 fwd), same at every $T$. \\
Cost of an $N$-site move & Same as single-site (ratio 0.998). \\
GPU memory per job & 6.4 GB allocated by PyTorch; 4.6--8.7 GB per process; 2 jobs per 24 GB card. \\
Does multi-site pay at equilibrium? & Open; pilot pending. Wins for relaxation, loses next to the reference. \\
Equilibrated sweep? & Not yet: 1000 moves is far from equilibrium at every $T\ge56$. \\
\bottomrule
\end{tabular}
\end{center}

\end{document}
